\documentclass[11pt,a4paper]{article}

% --- Encodage / langue ---
\usepackage[utf8]{inputenc}
\usepackage[T1]{fontenc}
\usepackage[french]{babel}
\usepackage{lmodern}

% --- Mise en page compacte ---
\usepackage[a4paper,margin=2cm]{geometry}
\usepackage{microtype}
\usepackage{enumitem}
\setlist{nosep,leftmargin=1.2em}
\usepackage{amsmath}
\usepackage{booktabs}
\usepackage{titlesec}
\titlespacing*{\section}{0pt}{8pt}{2pt}
\titlespacing*{\subsection}{0pt}{5pt}{1pt}
\titleformat{\section}{\normalfont\large\bfseries}{\thesection.}{0.5em}{}
\titleformat{\subsection}{\normalfont\bfseries}{}{0em}{}
\setlength{\parindent}{0pt}
\setlength{\parskip}{2pt}

\usepackage[backend=biber,style=numeric-comp,sorting=none,maxbibnames=3]{biblatex}
\addbibresource{references.bib}
\renewcommand*{\bibfont}{\footnotesize}

\usepackage{hyperref}
\hypersetup{hidelinks}

\newcommand{\TP}{\mathrm{TP}}
\newcommand{\FP}{\mathrm{FP}}
\newcommand{\FN}{\mathrm{FN}}

\title{\vspace{-1.2cm}\textbf{Livrable 2 --- Choix des méthodes d'évaluation\\[2pt]
\large Analyse automatique de partitions musicales (OMR)}}

\author{\normalsize Thibault Bourdoiseau \& Léo Li --- SAE TIFO 1 / TIFO 2 / MLRF}
\date{\vspace{-8pt}\small \today}

\begin{document}
\maketitle
\vspace{-10pt}

\section{Contexte}
Le système à évaluer est une chaîne de reconnaissance optique de partitions (OMR) qui transforme l'image
d'une partition imprimée --- voire manuscrite, pour éprouver sa robustesse --- en une liste de notes
(hauteur et durée), exportable au format MIDI~\cite{calvo2020,rebelo2012}. On évalue à \textbf{trois critères numériques complémentaires}, chacun
ciblant une étape de la chaîne~: la \emph{précision des hauteurs} et la \emph{précision des durées} mesurent
le décodage musical (le résultat final), tandis que le \emph{F1 de détection} mesure en amont la localisation
des symboles. Chaque critère compare une \textbf{prédiction} produite par le programme à une
\textbf{vérité-terrain}. Toutes les métriques sont vues en cours~: la F-mesure (précision/rappel) et la
distance d'édition sont dans la note~1 de MLRF, l'IoU (indice de Jaccard) en séance~6, et l'évaluation de
classifieurs en séances~4--5.

\section{Critère 1 --- Précision des hauteurs (Pitch Accuracy)}
\textbf{Comment se calcule-t-il (critère numérique).} On apparie d'abord la séquence de notes prédite à la
séquence de référence par \textbf{distance d'édition} (Levenshtein~\cite{levenshtein1966}), puis sur les
$N$ notes de référence~:
\[
\mathrm{PA} = \frac{1}{N}\sum_{i=1}^{N}\mathbf{1}\!\left[\hat h_i = h_i\right]\times 100,
\]
où $h_i$ est la hauteur de référence (p.~ex.\ \texttt{C4}) et $\hat h_i$ la hauteur prédite alignée (une note
non appariée compte comme fausse). Une note est correcte si \emph{nom et octave} coïncident. L'alignement par
édition (plutôt que position par position) évite qu'une seule note manquée décale toutes les suivantes et
effondre artificiellement le score.
\textbf{Ce qu'il caractérise.} La qualité du \emph{décodage vertical}~: lecture de la position de la tête de
note sur la portée, prise en compte de la clef et des altérations. C'est le critère le plus directement lié à
l'utilisabilité de l'export MIDI.
\textbf{Valeur raisonnable.} $\mathrm{PA}>90\,\%$ excellent (usage réel), $75\text{--}90\,\%$ correct
(prototype fonctionnel), $<75\,\%$ insuffisant (décodage à revoir). Valeurs dans $[0,100]\,\%$.

\section{Critère 2 --- Précision des durées (Rhythm Accuracy) et Note Accuracy}
\textbf{Comment se calcule-t-il.} Sur le même alignement que le critère~1~:
\[
\mathrm{RA} = \frac{1}{N}\sum_{i=1}^{N}\mathbf{1}\!\left[\hat d_i = d_i\right]\times 100,
\qquad
\mathrm{NA} = \frac{1}{N}\sum_{i=1}^{N}\mathbf{1}\!\left[\hat h_i = h_i\right]\mathbf{1}\!\left[\hat d_i = d_i\right]\times 100,
\]
où $d_i \in \{$ronde, blanche, noire, croche, double croche, \dots$\}$. La \textbf{Note Accuracy} exige
hauteur \emph{et} durée correctes simultanément~; on a toujours $\mathrm{NA}\le\min(\mathrm{PA},\mathrm{RA})$,
ce qui en fait un résumé sévère mais fidèle du résultat exploitable.
\textbf{Ce qu'il caractérise.} La \emph{classification des symboles de durée}~: type de tête (pleine ou vide),
présence d'une hampe, de ligatures et de points. La RA est souvent plus dure à atteindre que la PA car elle
dépend de plusieurs étapes successives.
\textbf{Valeur raisonnable.} Mêmes seuils que la PA ($>90\,\%$ usage réel, $>80\,\%$ prototype). Les erreurs
de durée sont particulièrement pénalisantes pour le MIDI car elles décalent toutes les notes suivantes dans
le temps. Valeurs dans $[0,100]\,\%$.

\section{Critère 3 --- F1 de détection des têtes de notes}
\textbf{Comment se calcule-t-il.} On compare les boîtes englobantes prédites aux boîtes de référence. Une
détection est un vrai positif ($\TP$) si son recouvrement $\mathrm{IoU}\ge\tau$ avec une boîte de référence
($\tau=0{,}5$, convention PASCAL VOC~\cite{everingham2010})~; l'appariement est glouton (une référence ne sert
qu'une fois). Les $\FP$ sont les détections sans correspondance, les $\FN$ les têtes non détectées~:
\[
\text{Précision}=\frac{\TP}{\TP+\FP},\qquad
\text{Rappel}=\frac{\TP}{\TP+\FN},\qquad
\text{F1}=\frac{2\cdot\text{Préc.}\cdot\text{Rappel}}{\text{Préc.}+\text{Rappel}}.
\]
\textbf{Ce qu'il caractérise.} L'étape de \emph{segmentation et de détection} des objets, indépendamment du
décodage musical. Il sert de diagnostic~: un F1 élevé avec une PA faible indique que les notes sont bien
localisées mais mal interprétées (clef, position), tandis qu'un F1 faible pointe la détection.
\textbf{Valeur raisonnable.} $\text{F1}>0{,}80$ attendu pour un prototype, $>0{,}90$ pour un système robuste.
On rapporte précision et rappel séparément, car ils orientent le réglage de la chaîne (seuils morphologiques,
filtrage des composantes connexes). Valeurs dans $[0,1]$.

\section{Outils de calcul}
Les trois critères sont implémentés dans \texttt{tools/metrics.py} (NumPy uniquement, aucune dépendance
lourde, pour faciliter le portage C++ du WP3)~: \texttt{detection\_f1} (critère~3) renvoie
précision/rappel/F1 à partir des boîtes~; \texttt{note\_scores} (critères~1--2) renvoie PA, RA et NA à partir
des deux séquences de notes (alignement par distance d'édition). Le script \texttt{tools/evaluate.py} lit une
vérité-terrain et une prédiction et affiche les scores. La base d'évaluation est décrite dans le document
séparé \emph{Livrable~2 --- Base d'évaluation}.

\printbibliography[heading=bibnumbered,title={Références}]

\end{document}
